\section{Admissible scaling at the current IUT interface and an exact order index}
\label{sec:iut-admissible-scaling-order-index}

We now compare the preceding source-faithful bridge with Project LANA commit
\texttt{c65b28c} \cite{LANAleanCurrent}.  This snapshot contains a genuine
positive advance: \path{Iut/Concrete/Existence.lean} constructs the previously
abstract \texttt{ConcreteThetaData\allowbreak Existence} datum from an elliptic
curve and its
arithmetic, Tate, representation-theoretic, and anabelian inputs.  The new
theorem \path{CurveInputs.concreteThetaDataExistence} therefore removes one
assembly gap.  The endpoint
\path{cor312Variant_implies_abc_curves}, however, remains conditional.  Its
hypotheses still include
\begin{itemize}
\item the height-theory proof package
  \path{Genl.HeightTheory.ProofPackage} and the curve inputs;
\item anabelian existence, together with providers for local theory,
  theta-local data, and tower arithmetic; and
\item the analytic prime bounds and the concrete Corollary~3.12 comparison.
\end{itemize}
It is therefore not an unconditional proof of Statement~I or of integer
$abc$.

The same snapshot retains a logically impossible quantifier in its local
log-volume interface.  Abstract the issue as follows.  Let $s:X\to X$, let
$V:\mathcal P(X)\to\R$, and fix $\delta\in\R$.

\begin{theorem}[Total-set and nonempty-set obstructions]
\label{thm:iut-total-set-scaling-no-go}
If
\begin{equation}\label{eq:iut-total-set-shift}
             V(s^{-1}U)=V(U)+\delta
             \qquad(U\subseteq X),
\end{equation}
then $\delta=0$.  If $X$ is nonempty, the same conclusion holds when
\eqref{eq:iut-total-set-shift} is required only for nonempty $U$.
More generally, the conclusion holds on any asserted domain containing a set
$U$ with $s^{-1}U=U$.
\end{theorem}

\begin{proof}
For the total law take $U=\varnothing$; its preimage is again empty and
additive cancellation gives $\delta=0$.  Under the nonempty restriction take
$U=X$, whose preimage is $X$.  The last assertion is the same cancellation
with the stated fixed set.
\end{proof}

In the current \texttt{LocalTheory.componentVol\_prime\_preimage} field one
has $\delta=\log p$ for a prime $p$.  Since $p\ge2$, this shift is positive,
so no inhabitant of the literal all-set structure can exist.  Merely appending
\texttt{U.Nonempty} would move the contradiction from the empty set to the
whole set and would not repair it.

The source-level repair used in our audit makes the domain explicit.  It adds
a predicate $\mathcal A$ of admissible finite, nonzero-volume regions and
requires
\begin{equation}\label{eq:iut-admissible-shift}
 U\in\mathcal A\Longrightarrow U\ne\varnothing,\qquad
 U\in\mathcal A\Longrightarrow s^{-1}U\in\mathcal A,\qquad
 U\in\mathcal A\Longrightarrow V(s^{-1}U)=V(U)+\log p.
\end{equation}
The predicate is transported from \texttt{LocalTheory} to
\texttt{LogVolumeData} and through \path{Concrete.Container.vol}.  No
numeric equality is imposed on sets outside its domain.  This is the minimal
quantifier correction suggested by the finite, nonzero-volume regime of
IUT~III, Proposition~3.9 and Remark~3.9.5 \cite{IUT3}.

\begin{proposition}[Admissible orbit law]
\label{prop:iut-admissible-orbit-law}
Assume \eqref{eq:iut-admissible-shift}.  For $U\in\mathcal A$ and $r\ge0$,
\[
                    V(s^{-r}U)=V(U)+r\log p.
\]
If $\log p\ne0$, the admissible preimage orbit of $U$ is injective and in
particular is not periodic.
\end{proposition}

\begin{proof}
Induction on $r$ proves the formula: closure under preimage supplies the
admissibility hypothesis needed at each successor step.  If
$s^{-r}U=s^{-t}U$, the formula gives
$(r-t)\log p=0$, whence $r=t$.
\end{proof}

The patched versions of the three affected source files compile together
with both upstream targets \texttt{Iut} and \texttt{Iut4Sec1}; the build
finishes all $8767$ tasks.  This verifies that the repaired dependency shape
is compatible with the current source tree.  It does not construct the local
Haar volume or prove the comparison theorem.

There is a second accounting seam at the passage from an integral tensor
order to the product of maximal factor orders.  For $n\in\N$ put
\begin{equation}\label{eq:iut-congruence-order}
 \mathcal O_n=\{(a,b)\in\Z^2:n\mid a-b\},\qquad
 \delta_n(a,b)=a-b\pmod n.
\end{equation}

\begin{theorem}[Exact congruence-order quotient]
\label{thm:iut-congruence-order-index}
The map $\delta_n:\Z^2\to\Z/n\Z$ is a surjective additive homomorphism with
kernel $\mathcal O_n$.  Consequently
\[
       \Z^2/\mathcal O_n\simeq\Z/n\Z,
       \qquad [\Z^2:\mathcal O_n]=n.
\]
For $n>1$, $\mathcal O_n$ is a proper finite-index subring of $\Z^2$, although
both coordinate projections $\mathcal O_n\to\Z$ are surjective.
\end{theorem}

\begin{proof}
Additivity is immediate, and every residue class is the image of $(a,0)$.
Moreover $\delta_n(a,b)=0$ exactly when $n\mid a-b$, so the first isomorphism
theorem gives the quotient.  The kernel is a subring: it contains $(0,0)$
and $(1,1)$, is additively closed, and for two of its elements
\[
                    ac-bd=a(c-d)+d(a-b),
\]
so it is closed under componentwise multiplication.  Its quotient
cardinality is $n$.  If $n>1$, the point
$(1,0)$ is absent, proving properness; every $(z,z)$ is present, proving both
projection statements.
\end{proof}

The theorem includes $n=0$ under the \texttt{Nat.card} index convention:
both the quotient and $\Z/0\Z\cong\Z$ are infinite and their recorded natural
cardinality is zero.  For positive $m,n$, exactness also gives
\begin{equation}\label{eq:iut-order-index-log-additivity}
 [\Z^2:\mathcal O_{mn}]
 =[\Z^2:\mathcal O_m][\Z^2:\mathcal O_n],
 \qquad
 \log[\Z^2:\mathcal O_{mn}]
 =\log[\Z^2:\mathcal O_m]+\log[\Z^2:\mathcal O_n].
\end{equation}
Thus a finite integral-order defect contributes an additive logarithmic term.
Identifying the actual IUT tensor-order quotient with a local different or
conductor exponent is a further arithmetic theorem, not a consequence of
this model.

The empty and whole sets are full-premise counterexamples to the two stated
overbroad volume laws.  The orders $\mathcal O_n$, $n>1$, are full-premise
counterexamples to the assertion that surjective component projections force
the product order.  These witnesses retire precisely those shortcuts.  They
do not satisfy the premises of an admissible, index-aware same-pilot theorem,
so that corrected route remains active.  Its unresolved obligations are to
construct the actual admissible Haar regions, calculate the tensor-order
indices at every factor and place, prove covariance through the horizontal
link and Ind3, and establish the concrete Corollary~3.12 comparison.

The Lean companion proves the two no-go theorems, identifies the exact kernel
and quotient in \eqref{eq:iut-congruence-order}, proves the index and its
logarithmic multiplication law, and verifies properness together with both
surjective projections.  The source ledger separately freezes the upstream
commit, exact patch replay, and complete build log.  No same-pilot theorem,
Statement~I, or $abc$ consequence is introduced as a formal assumption.
